\documentclass[11pt]{article}
\usepackage[T1]{fontenc}
\usepackage[utf8]{inputenc}
\usepackage{lmodern,amsmath,amssymb,amsthm,mathtools}
\usepackage[margin=25mm]{geometry}
\usepackage[hidelinks]{hyperref}
\newcommand{\C}{\mathbb C}
\newcommand{\R}{\mathbb R}
\newcommand{\PP}{\mathcal P}
\newcommand{\Bcal}{\mathcal B}
\newcommand{\Rcal}{\mathcal R}
\newcommand{\Scal}{\mathcal S}
\newcommand{\norm}[1]{\left\lVert#1\right\rVert}
\newcommand{\ip}[2]{\left\langle#1,#2\right\rangle}
\DeclareMathOperator{\Tr}{Tr}
\DeclareMathOperator{\diag}{diag}
\DeclareMathOperator{\im}{im}
\DeclareMathOperator{\rem}{rem}
\DeclareMathOperator{\rad}{rad}
\title{Central-path transfer for the original arithmetic annihilator metrics}
\author{Split-Zero research continuation}
\date{19 September 2026}
\begin{document}
\maketitle

\begin{abstract}
The full two-parameter arithmetic source interpolation has a uniformly
small central derivative on every specified complete annihilator
polynomial space.  Its explicit derivative bound is the product of the
original logarithmic density contrast and the central localization
fraction.  Integrating this derivative proves exponential relative
metric bounds and an explicit bound for the actual minimum sections.
For the complete relation ideal this gives a finite central-leg error
and a finite rectangle error in the original four-endpoint signed
determinant.  The remaining source-band interaction is expressed by its
exact two-variable reproducing-kernel integral.  Every statement retains
the full primary orders, original physical coordinate, original
polynomial remainder map, complete real integration domain and both
endpoint multiplicities.
\end{abstract}

\section{Original source, analytic providers and fixed notation}

Fix the original complete hypothetical quartet
\[
 \rho=\tfrac12+\delta+i\gamma,\qquad
 0<\delta<\tfrac12,\quad\gamma>2,\quad m_0\in\mathbb Z_{\ge1}.
\]
The divisor, its entire Taylor unit, and its arithmetic density are
\begin{equation}\tag{CPT1}
\begin{gathered}
 h(z)=\prod_{\lambda\in\{\rho,\bar\rho,1-\rho,1-\bar\rho\}}
             (z-\lambda)^{m_0},\qquad
 v_h(z)=2\xi(z)/h(z),\\
 w_h(y)=|v_h(1/2+iy)|^2/(2\pi),\qquad m_k=w_h^{*k},\\
 k\equiv1\pmod4,\quad k\ge5,\quad
 e=1+k(m_0-1),\quad q=e(k+1)^2,\quad c=k/2,\\
 \chi(S)=\prod_{a,b=0}^k
  [S-c-(2a-k)\delta-i(2b-k)\gamma]^e,\qquad
 E=\C[S]/(\chi).
\end{gathered}
\end{equation}
The symbol \(\xi\) and the zero hypothesis have their original
Riemann-xi meaning.  This note computes consequences on that original
counterfactual source; it supplies no numerical off-critical zero.
Every limit below fixes the complete packet.

The Gamma measure, full tilted mass and central endpoint are
\begin{equation}\tag{CPT2}
\begin{gathered}
 \sigma(y)=\frac{|\Gamma(1/4+iy/2)|^2}{2\pi},\quad
 d\sigma=\sigma(y)\,dy,\quad c_\sigma=\int_{\R}d\sigma=\sqrt{2\pi},
 \quad \alpha=\pi/2,\\
 M_h(\theta)=\int_{\R}e^{\theta y}w_h(y)\,dy,\quad
 \ell_h(\theta)=\log M_h(\theta),\quad
 M=M_h(\alpha),\quad M_0=M_h(0),\\
 \mu=\ell_h'(\alpha)>0,\quad T=k\mu,\quad
 R=k\sqrt{\delta^2+\gamma^2},\quad
 C=[-T,T],\quad O=\R\setminus C.
\end{gathered}
\end{equation}
The notation \(M\) here always denotes the full tilted mass; set
\(\beta_0=21+4m_0\) for the polynomial exponent, to avoid confusing
these two different quantities.  In particular \(B_0=2\beta_0=42+8m_0\)
and \(p=8m_0-9/2>3\).

The following original analytic estimates are the GTM1 and RTM2
providers, including their constants \cite{providers}:
\begin{equation}\tag{CPT3}
\begin{gathered}
 a_k(1+y^2)^{-\beta_0}\sigma(y)\le m_k(y)\le A_k\sigma(y),\\
 a_k=
 \frac{c_h\vartheta_h^{\,k-3}e^{-\alpha(k-3)}(k-2)^{-B_0}}
 {C_\Gamma2^{B_0/2}},\qquad
 A_k=\frac{C_h^{\rm tilt}}{c_\Gamma}k^{p+1}M^{k-1},
 \qquad \vartheta_h=\int_{-1}^1w_h(y)\,dy,\\
 c_\Gamma e^{-\alpha|y|}(1+|y|)^{-1/2}
 \le\sigma(y)\le
 C_\Gamma e^{-\alpha|y|}(1+|y|)^{-1/2},\\
 c_\Gamma=\sqrt2e^{-7/6},\qquad
 C_\Gamma=\sqrt{2\sqrt5}e^{2/3}.
\end{gathered}
\end{equation}
Here \(c_h,C_h^{\rm tilt}>0\) are the original threefold convolution
and boundary-tilt constants.  Their original source bounds are
\[
 w_h^{*3}(y)\ge c_h e^{-\alpha|y|}(1+|y|)^{-B_0},
 \qquad e^{\alpha y}w_h(y)\le
 C_h^{\rm tilt}(1+|y|)^{-p}.
\]
These statements, with their source derivations, are preserved in the
provider file; this continuation does not select new density constants.
No positive lower estimate for the one-factor zeta numerator is used.

Define the exact logarithmic contrast and the entire parameter square:
\begin{equation}\tag{CPT4}
\begin{gathered}
 H(y)=\log\frac{m_k(y)}{M^k\sigma(y)},\qquad
 H_C=1_CH,\quad H_O=1_OH,\\
 d\nu_{s,t}(y)=\sigma(y)e^{sH_C(y)+tH_O(y)}\,dy,
 \qquad (s,t)\in[0,1]^2.
\end{gathered}
\end{equation}
Thus \(\nu_{0,0}=\sigma\), \(\nu_{1,1}=M^{-k}m_k(y)dy\), and
\(\nu_{0,1}\) retains the exact arithmetic exterior while its centre is
the specified Gamma density.  The factor \(M^k\) has not disappeared:
its exact cancellation in the particular signed determinant is proved
in Section~\ref{sec:determinants}.

All polynomial inner products are conjugate linear in their first
argument:
\[
 \ip{P}{Q}_{s,t}
 =\int_{\R}\overline{P(c+iy)}Q(c+iy)\,d\nu_{s,t}(y).
\]
The coefficient map
\begin{equation}\tag{CPT5}
 U_{c,D}:\C[S]_{\le D}\longrightarrow\C[y]_{\le D},
 \qquad U_{c,D}P(y)=P(c+iy)
\end{equation}
is an exact complex-linear isomorphism with inverse
\(F(y)\mapsto F((S-c)/i)\).  In increasing monomial frames its
diagonal is \(1,i,\ldots,i^D\), so
\(\det U_{c,D}=i^{D(D+1)/2}\).  This coordinate map, with modulus-one
determinant, will be displayed whenever the real-coordinate proof is
returned to the original coefficient quotient.

\section{Uniform bounds on the entire interpolation square}

Put
\begin{equation}\tag{CPT6}
 b_*=\max\{1,A_k/M^k\},\quad a_*=\min\{1,a_k/M^k\},
 \qquad
 a_D=\frac{a_*}{[1+4(D+1)^2]^{\beta_0}}.
\end{equation}
For every \((s,t)\in[0,1]^2\),
\begin{equation}\tag{CPT7}
 a_*(1+y^2)^{-\beta_0}\,d\sigma
 \le d\nu_{s,t}\le b_*\,d\sigma,\qquad
 a_D\norm{F}_\sigma^2\le\norm{F}_{s,t}^2
 \le b_*\norm{F}_\sigma^2\quad(\deg F\le D).
\end{equation}

To prove the pointwise bounds, let
\(r(y)=m_k(y)/(M^k\sigma(y))>0\).
At each \(y\), precisely one exponent in (CPT4) acts, and its value
is some \(u\in[0,1]\).  For any positive \(r\),
\[
 \min\{1,r\}\le r^u\le\max\{1,r\}.
\]
Applying (CPT3) gives the first two bounds in (CPT7).
In particular this proves a pointwise bound before any polynomial norm
is considered.

For the polynomial lower estimate, the monic Gamma polynomials
\(b_j(y)\) obey
\begin{equation}\tag{CPT8}
 yb_j=b_{j+1}+j(j-\tfrac12)b_{j-1},\qquad
 \gamma_j=\norm{b_j}_\sigma^2
 =\sqrt{2\pi}\frac{(2j)!}{4^j}.
\end{equation}
These are GTM7 \cite{providers}; the recurrence is also the
Meixner--Pollaczek recurrence in the convention
\(b_j(y)=j!P_j^{(1/4)}(y/2;\pi/2)\)
\cite[(18.22.8)]{dlmfMP}.  Indeed substituting
\(\lambda=1/4,\varphi=\pi/2,x=y/2\) into that cited recurrence and
multiplying by \(j!\) gives exactly (CPT8).
For the orthonormal polynomials \(e_j=b_j/\sqrt{\gamma_j}\),
\[
 ye_j=\sqrt{(j+1)(j+\tfrac12)}\,e_{j+1}
       +\sqrt{j(j-\tfrac12)}\,e_{j-1}.
\]
The upward and downward shifts on
\(\operatorname{span}(e_0,\ldots,e_D)\) each have operator norm at
most \(D+1\), including the output coefficient at degree \(D+1\).
The triangle inequality therefore proves
\[
 \norm{yF}_\sigma\le2(D+1)\norm{F}_\sigma.
\]
For a nonzero \(F\) divide the positive measure
\(|F|^2d\sigma\) by its actual mass \(\norm{F}_\sigma^2\), solely
to apply the convex scalar inequality
\(\mathbb E(1+X)^{-\beta_0}\ge(1+\mathbb EX)^{-\beta_0}\)
to \(X=y^2\).  This inequality follows from the supporting tangent
line to the function \((1+x)^{-\beta_0}\), whose second derivative
is positive on \(x\ge0\).  Multiplying the mass back gives
\[
 \int |F|^2(1+y^2)^{-\beta_0}d\sigma
 \ge[1+4(D+1)^2]^{-\beta_0}\norm{F}_\sigma^2.
\]
Combining this with the pointwise lower bound proves (CPT7).
No source measure or source polynomial is replaced in that argument.

The pointwise envelopes also imply
\[
 |H(y)|\le
 |\log(a_k/M^k)|+|\log(A_k/M^k)|
       +\beta_0\log(1+y^2).
\]
Thus every finite polynomial moment and every finite parameter
derivative of it is dominated on the closed square by a polynomial
times a power of \(\log(1+y^2)\) times the exponentially decreasing
Gamma density.  Differentiation under the integral is valid to all
finite orders, with one-sided derivatives at the square's boundary.
Every Gram is positive definite, since its density is positive
everywhere and a nonzero polynomial has only finitely many real
zeros.  Its inverse and every fixed-fibre minimum section therefore
depend smoothly on the square.

\section{Finite central localization with every root retained}
\label{sec:localization}

Let \(g\mid\chi\) be an actual monic divisor, with its selected full
primary orders, and write
\begin{equation}\tag{CPT9}
 J=\deg g,\quad d_0=q-J,\quad
 Q_g(y)=i^{-d_0}(\chi/g)(c+iy),\quad
 L=D-d_0,\quad
 V_{g,D}=Q_g\,\C[y]_{\le L}.
\end{equation}
Assume the explicit finite guards
\[
 0\le J\le q,\qquad d_0\le D\le2q+1,\qquad
 R<q,\qquad T\le q/2.
\]
Each root of \(Q_g\), counted with its remaining multiplicity, has
modulus at most \(R\).  Define the finite positive bound and its cap
\begin{equation}\tag{CPT10}
\begin{split}
 \mathfrak e_{k,J,D}
 ={}&\frac{b_*}{a_D}\frac{c_\sigma}{c_\Gamma}
 \frac{(L+1)^2\sqrt{1+2q}}q\,e^{\pi q}100^L
 \left(\frac{T+R}{q-R}\right)^{2(q-J)},\\
 &\varepsilon_{k,J,D}=\min\{1,\mathfrak e_{k,J,D}\}.
\end{split}
\end{equation}
Then the entire parameter square satisfies
\begin{equation}\tag{CPT11}
 \int_C |F(y)|^2d\nu_{s,t}(y)
 \le\varepsilon_{k,J,D}\norm{F}_{s,t}^2,
 \qquad F\in V_{g,D}.
\end{equation}

Here is the full polynomial estimate.  On \([q,2q]\) the
orthonormal Legendre basis of polynomials of degree at most \(L\)
is
\[
 \varphi_j(y)=q^{-1/2}\sqrt{2j+1}\,
 P_j(2y/q-3),\qquad 0\le j\le L.
\]
The Legendre norm \(2/(2j+1)\) and recurrence used here are
respectively \cite[Table 18.3.1]{dlmfLegendreNorm} and
\cite[Table 18.9.1]{dlmfLegendreRec}.  The latter gives
\[
 (j+1)P_{j+1}(x)=(2j+1)xP_j(x)-jP_{j-1}(x).
\]
For \(|x|\le4\), \(P_0=1,P_1=x\) and induction prove
\(|P_j(x)|\le10^j\): the inductive upper bound is
\[
 \frac{2j+1}{j+1}\,4\,10^j+
 \frac{j}{j+1}\,10^{j-1}
 \le8\,10^j+10^{j-1}<10^{j+1}.
\]
If \(|y|\le T\le q/2\), the argument \(2y/q-3\) belongs to
\([-4,-2]\).  Expand the full polynomial
\(P=\sum_{j=0}^L d_j\varphi_j\), without omitting any coefficient.
Orthogonality and Cauchy--Schwarz give
\begin{equation}\tag{CPT12}
 |P(y)|^2\le
 \left(\sum_{j=0}^L|d_j|^2\right)
 \left(\frac1q\sum_{j=0}^L(2j+1)100^j\right)
 \le\frac{(L+1)^2 100^L}{q}
                       \int_q^{2q}|P(x)|^2dx .
\end{equation}
For \(F=Q_gP\), all its roots give
\[
 |Q_g(y)|\le(T+R)^{d_0}\quad(y\in C),\qquad
 |Q_g(x)|\ge(q-R)^{d_0}\quad(q\le x\le2q).
\]
The upper bound in (CPT7), (CPT12) and the unchanged total Gamma
mass yield
\[
 \int_C|F|^2d\nu_{s,t}
 \le b_*c_\sigma(T+R)^{2d_0}
       \frac{(L+1)^2 100^L}{q}\int_q^{2q}|P(x)|^2dx.
\]
The polynomial lower bound in (CPT7) is applied on the \emph{whole}
real line before its Gamma integral is restricted to \([q,2q]\).
The exact Gamma envelope in (CPT3) gives
\[
 \norm{F}_{s,t}^2
 \ge a_D\frac{c_\Gamma e^{-\pi q}}{\sqrt{1+2q}}
           (q-R)^{2d_0}\int_q^{2q}|P(x)|^2dx .
\]
For \(P\ne0\), division gives \(\mathfrak e_{k,J,D}\).
The trivial central-fraction bound is one.  The zero polynomial
satisfies (CPT11) directly.  This proves (CPT11) with both full
real integration regions present in the norm.

For fixed complete packet and \(D\le2q+1\),
\(\log(b_*/a_D)=O_h(k+\log q)\).  Since \(T+R\) is \(k\) times
a fixed positive constant and \(R/q\to0\), direct logarithms in
(CPT10) give, for its uncapped positive bound,
\begin{equation}\tag{CPT13}
 \log\mathfrak e_{k,J,2q+1}
 =-2(q-J)\log(q/k)+O_h(q+J+k+\log q).
\end{equation}
For every fixed \(0<\vartheta<1\) and \(A>0\), this implies
\[
 q^A\mathfrak e_{k,J,2q+1}\longrightarrow0
 \quad\hbox{uniformly for }0\le J\le(1-\vartheta)q.
\]
The finite formula (CPT10), rather than this asymptotic statement,
is used in all bounds below.  At \(J=q\) the root factor is absent;
the capped bound remains valid but supplies no such decay.

\section{An explicit logarithmic central contrast}

The original central-source calculation ACW11--13
\cite{centralProvider} proves, on its explicit local-density guard
\(\epsilon_k^{\rm loc}\le1/2\), the following formula.  Let
\(\theta(u)\) invert \(\ell_h'\) on \([-\alpha,\alpha]\), and put
\[
 V(\theta)=\ell_h''(\theta),\qquad
 0<v_*\le V(\theta)\le V_*,
 \qquad c_{\sigma}^{\rm env}(y)
 =\sigma(y)e^{\alpha|y|}\sqrt{1+|y|}.
\]
Then, for \(|u|\le\mu\),
\begin{equation}\tag{CPT14}
\begin{split}
 H(ku)={}&kF_h(u)+\tfrac12\log(|u|+k^{-1})
 -\tfrac12\log(2\pi V(\theta(u)))
 -\log c_{\sigma}^{\rm env}(ku)
 +\log(1+\epsilon_{k,\theta(u)}),\\
 F_h(u)={}&
 \ell_h(\theta(u))-u\theta(u)-\ell_h(\alpha)
                           +\alpha|u|,\qquad
 |\epsilon_{k,\theta}|\le\epsilon_k^{\rm loc}.
\end{split}
\end{equation}
For completeness, all constants and the local-density proof yielding
this provider are spelled out in the appendix.  In particular this
formula retains the actual error, not a selected saddle coefficient.

The function \(F_h\) is even, and on \(0<u<\mu\),
\[
 F_h'(u)=\alpha-\theta(u)>0,\qquad F_h(\mu)=0,\qquad
 F_h(0)=\log M_0-\log M.
\]
Consequently \(-(\log M-\log M_0)\le F_h(u)\le0\).
The exact argument \(|u|+1/k\) lies between \(1/k\) and \(\mu+1\).
Set
\begin{equation}\tag{CPT15}
\begin{split}
 L_k={}&k(\log M-\log M_0)
 +\tfrac12\max\{\log k,|\log(\mu+1)|\}\\
 &+\tfrac12\max\{|\log(2\pi v_*)|,|\log(2\pi V_*)|\}\\
 &+\max\{|\log c_\Gamma|,|\log C_\Gamma|\}
 +2\epsilon_k^{\rm loc}.
\end{split}
\end{equation}
Since \(|\log(1+x)|\le2|x|\) for \(|x|\le1/2\),
(CPT14) proves the finite estimate
\begin{equation}\tag{CPT16}
 |H_C(y)|\le L_k,\qquad L_k=O_h(k+\log k).
\end{equation}
All terms in (CPT15) are nonnegative.  In particular
\(\log M-\log M_0\ge0\) by the original evenness of \(w_h\).

There is also a uniform upper bound on the positive part:
\begin{equation}\tag{CPT17}
 H(y)\le K_h\quad(y\in C),\qquad
 K_h=\left[
 \tfrac12\log(\mu+1)-\tfrac12\log(2\pi v_*)
 -\log c_\Gamma+\log2\right]_+.
\end{equation}
Indeed the \(kF_h\) term is nonpositive, and each other term of
(CPT14) has the corresponding upper bound in (CPT17).

\section{Central derivative, transported metrics and attained sections}
\label{sec:sections}

Fix \(g,D\) as in Section~\ref{sec:localization}, and abbreviate
\begin{equation}\tag{CPT18}
 \Lambda_{k,J,D}=L_k\varepsilon_{k,J,D}.
\end{equation}
On \(V_{g,D}\) the actual derivative form is
\[
 \dot h^C_{s,t}(F,G)
 =\partial_s\ip{F}{G}_{s,t}
 =\int_C H(y)\overline{F(y)}G(y)d\nu_{s,t}(y).
\]
The bound (CPT16), Cauchy--Schwarz on \(C\), and (CPT11) prove
\begin{equation}\tag{CPT19}
 |\dot h^C_{s,t}(F,G)|
 \le\Lambda_{k,J,D}\norm{F}_{s,t}\norm{G}_{s,t}.
\end{equation}
In any fixed coefficient frame, write \(A_{s,t}\) for the
positive Gram and \(A^C_{s,t}=\partial_s A_{s,t}\).
Thus (CPT19) is the proved operator inequality
\[
 \norm{A_{s,t}^{-1/2}A^C_{s,t}A_{s,t}^{-1/2}}_{\rm op}
 \le\Lambda_{k,J,D}.
\]
For a fixed nonzero \(F\), its squared norm satisfies
\(|\partial_s\log\norm{F}_{s,t}^2|\le\Lambda_{k,J,D}\).
Integrating from \(s_0\) to \(s_1\) gives
\begin{equation}\tag{CPT20}
 e^{-\Lambda_{k,J,D}|s_1-s_0|}A_{s_0,t}
 \preceq A_{s_1,t}
 \preceq e^{\Lambda_{k,J,D}|s_1-s_0|}A_{s_0,t}.
\end{equation}
This bound holds for every \(t\in[0,1]\).
It requires no smallness assumption on
\((e^{L_k}-1)\varepsilon_{k,J,D}\).

Here is the exact transfer through every fixed restriction and
attained quotient of this polynomial space.  A fixed injective
matrix \(I\) gives the restriction \(I^*A_{s,t}I\); congruence
preserves (CPT20).  A fixed onto matrix
\(J:V_{g,D}\twoheadrightarrow Y\), with a fixed value frame in \(Y\),
has minimum metric and section
\begin{equation}\tag{CPT21}
 G^Y_{s,t}=(J A_{s,t}^{-1}J^*)^{-1},\qquad
 R^Y_{s,t}=A_{s,t}^{-1}J^*G^Y_{s,t}.
\end{equation}
Direct multiplication gives \(JR^Y_{s,t}=I_Y\).
For \(z\in\ker J\),
\[
 z^*A_{s,t}R^Y_{s,t}v=z^*J^*G^Y_{s,t}v=0.
\]
Every representative of \(v\) is \(R^Y_{s,t}v+z\); orthogonality
proves that its squared norm is
\(v^*G^Y_{s,t}v+z^*A_{s,t}z\).
This proves both formulas in (CPT21) and attainment of the unique
minimum.  Taking the infimum over the \emph{same} fibre in each
side of (CPT20) gives the same bounds for \(G^Y\).
Repeating these two operations proves the bounds for any stated
fixed sequence of restrictions and attained quotients inside this
ambient polynomial space.  For a rank-\(r\) resulting metric \(G\),
\begin{equation}\tag{CPT22}
 \left|\log\frac{\det G_{s_1,t}}{\det G_{s_0,t}}\right|
 \le r\Lambda_{k,J,D}|s_1-s_0|.
\end{equation}
The rank is retained, including rank zero with determinant one.

There is a direct estimate for the actual minimizing polynomials.
Let \(P_s=R^Y_{s,t}v\), for \(v\in Y\) fixed and \(t\) fixed.
Because \(JP_s=v\), differentiation gives \(P_s'\in\ker J\).
For every \emph{fixed} \(z\in\ker J\), differentiate its exact
orthogonality to obtain
\[
 \ip{z}{P_s'}_{s,t}=-\dot h^C_{s,t}(z,P_s).
\]
This identity holds for every vector of the fixed kernel at each
parameter, so it may then be evaluated at \(z=P_s'\).
Using (CPT19) and dividing when \(P_s'\ne0\) proves
\begin{equation}\tag{CPT23}
 \norm{P_s'}_{s,t}\le
 \Lambda_{k,J,D}\norm{P_s}_{s,t}.
\end{equation}
The minimum property and (CPT20) give
\(\norm{P_s}_{s,t}\le e^{s\Lambda_{k,J,D}/2}
\norm{P_0}_{0,t}\).
The same metric comparison gives
\(\norm{P_s'}_{0,t}\le e^{s\Lambda_{k,J,D}/2}
\norm{P_s'}_{s,t}\).
Integrating the resulting bound
\(\norm{P_s'}_{0,t}\le
\Lambda_{k,J,D}e^{s\Lambda_{k,J,D}}\norm{P_0}_{0,t}\)
in the fixed coefficient space proves
\begin{equation}\tag{CPT24}
 \boxed{\norm{P_1-P_0}_{0,t}
 \le(e^{\Lambda_{k,J,D}}-1)\norm{P_0}_{0,t}.}
\end{equation}
At \(\Lambda=0\) this formula follows by the zero derivative.
For arbitrary endpoints the same proof, with the oriented segment
reparametrized, gives \(e^{\Lambda|s_1-s_0|}-1\) and the starting
norm.  Thus the original coefficients and phases of the entire
minimum section, rather than only its determinant, are controlled.

\section{The exact map to the original primary submodule}

Let \(A:E\to E\) denote multiplication by the original \(S\).
The polynomial-domain division identity gives
\begin{equation}\tag{CPT25}
 E[g]:=\ker g(A)=(\chi/g)\C[S]/(\chi),\qquad
 \dim_\C E[g]=J.
\end{equation}
Indeed \(g[P]_\chi=0\) says that \(\chi\mid gP\).
Writing \(\chi=g(\chi/g)\) and cancelling the nonzero polynomial
\(g\) in \(\C[S]\) is equivalent to \((\chi/g)\mid P\).
Multiplication by \(\chi/g\) identifies
\(\C[S]/(g)\) bijectively with this submodule, proving the
dimension assertion and retaining every primary multiplicity.

For \(N\ge q-1\), put
\[
 W_{g,N}=(\chi/g)\C[S]_{\le N-q+J},\qquad
 J_{g,N}:W_{g,N}\longrightarrow E[g],\quad
 J_{g,N}P=[P]_\chi.
\]
For \(J=0\) the value space is zero and the formula has that
literal meaning.  The onto property follows from the
degree-\(<J\) representatives in \(\C[S]/(g)\).
Its kernel is exactly \(\chi\C[S]_{\le N-q}\), with negative
degree interpreted as the zero space.  Moreover every
degree-\(\le N\) representative in the original remainder fibre
of a class in \(E[g]\) already lies in \(W_{g,N}\), by the
divisibility just proved.  Therefore minimizing in \(W_{g,N}\)
is precisely the original minimum for that class.

In real coordinates \(U_{c,N}W_{g,N}=V_{g,N}\), because the
constant factor \(i^{q-J}\) changes no subspace.  This gives the
commuting diagram in explicit map form:
\begin{equation}\tag{CPT26}
\begin{gathered}
 W_{g,N}\ \xrightarrow{\ U_{c,N}\ }\ V_{g,N},\\
 J_{g,N}
 =\left(P\longmapsto[P]_\chi\right)\big|_{W_{g,N}},
 \qquad
 J_{g,N}\circ U_{c,N}^{-1}:V_{g,N}\twoheadrightarrow E[g].
\end{gathered}
\end{equation}
If \(G_N[\nu_{s,t}]\) is the original full quotient Gram on \(E\)
and \(I_g\) is a fixed frame inclusion of \(E[g]\) into \(E\),
the minimum metric in (CPT21) is exactly
\begin{equation}\tag{CPT27}
 G^{g}_{N;s,t}=I_g^*G_N[\nu_{s,t}]I_g.
\end{equation}
Equations (CPT20)--(CPT24) consequently apply to the original
restriction, its minimum sections and its fixed subquotients,
with \(D=N\), or with \(D=2q+1\) as one common ambient degree.
In particular, at all four original cutoff degrees the same
explicit large-degree constant is valid.

An opposite terminal eigenpair uses the actual divisor
\(g=(S-\lambda)(S-(k-\lambda))\) and \(J=2\), including when
the full primary length \(e\) exceeds one.  The remaining
\(e-1\) powers at these roots stay in \(\chi/g\).
All terminal eigenlines together use \(g=\rad\chi\) and
\(J=(k+1)^2=q/e\).  For \(m_0>1\) this whole terminal space
has the decay (CPT13).  For \(m_0=1\), this last divisor is
\(\chi\), hence \(J=q\); it is the full quotient.

The exact determinant relation records the complementary metric
that is still present when \(J<q\).  Complete \(I_g\) by any fixed
frame \(I_\perp\) to a frame of \(E\), and in that frame write
\[
 G_N=\begin{pmatrix} A_g&B_g\\B_g^*&C_g\end{pmatrix},
 \qquad A_g=I_g^*G_NI_g.
\]
The identity induced on the quotient \(E/E[g]\) has attained
metric
\begin{equation}\tag{CPT28}
 Q_g^{\rm comp}=C_g-B_g^*A_g^{-1}B_g,\qquad
 \det G_N=\det A_g\,\det Q_g^{\rm comp}.
\end{equation}
To verify it, a value \(z\) in this quotient is represented by
\((x,z)\).  Completing the square in its norm gives the unique
minimum at \(x=-A_g^{-1}B_gz\) and the displayed \(Q_g^{\rm comp}\).
The block triangular elimination has determinant one, proving
the determinant formula.  Thus the map \(E[g]\hookrightarrow
E\twoheadrightarrow E/E[g]\) and its precise complementary
metric remain in the calculation.  A bound on \(A_g\) alone
has not evaluated this second factor.

\section{Original Schur identity and its signed central receiver}
\label{sec:determinants}

For any source \(\nu_{s,t}\) define the original degree-\(N\)
moment matrix in the \(S\)-monomial frame by
\[
 (H_N)_{ij}=\int_\R
       \overline{(c+iy)^i}(c+iy)^j\,d\nu_{s,t}(y),
 \qquad 0\le i,j\le N.
\]
Let \(\omega_n(s,t)\) be its squared monic source minimum of
degree \(n\), and let \(\eta_r(s,t)\) be the squared monic
relation minimum
\[
 \eta_r(s,t)=
 \min_{\substack{u\ {\rm monic}\\\deg u=r}}
 \int_\R|\chi(c+iy)u(c+iy)|^2\,d\nu_{s,t}(y).
\]
Every preceding source polynomial, respectively preceding
relation \(\chi\C[S]_{<r}\), is included in these minima.
Gram--Schmidt in monic triangular frames gives
\[
 \det H_N=\prod_{n=0}^N\omega_n,\qquad
 \det L_r=\prod_{j=0}^{r-1}\eta_j,
 \qquad
 L_r=\operatorname{Gram}(\chi,\chi S,\ldots,\chi S^{r-1}),
 \quad \det L_0=1.
\]
These are identities of the actual source moments.

The concatenated coefficient frame
\[
 1,S,\ldots,S^{q-1},\chi,\chi S,\ldots,\chi S^{N-q}
\]
is monic triangular, with determinant one.
Its lower Gram block is \(L_{N-q+1}\).  Completing the square
in its relation coefficients yields precisely the remainder
minimum metric
\[
 G_N=(J_NH_N^{-1}J_N^*)^{-1},\qquad
 J_NP=[P]_\chi.
\]
Taking the block determinant proves
\begin{equation}\tag{CPT29}
 \det G_N=
 \frac{\prod_{n=0}^N\omega_n}
      {\prod_{j=0}^{N-q}\eta_j}\quad(N\ge q-1).
\end{equation}
For \(N=q-1\) the denominator is one.
Under (CPT5) a monic degree-\(n\) polynomial becomes one of
leading coefficient \(i^n\); multiplying it by \(i^{-n}\)
returns monicity and leaves its norm unchanged.  The relation
\(\chi u\) has the corresponding factor \(i^{q+r}\).
Hence no determinant modulus or monic norm acquires a phase
factor in (CPT29).

Retain the four signs, the full relation band and its weights:
\begin{equation}\tag{CPT30}
\begin{gathered}
 \Bcal(s,t)=\log\det G_{q-1}+\log\det G_q
             -\log\det G_{2q-1}-\log\det G_{2q},\\
 c_0=c_q=1,\qquad c_r=2\quad(1\le r\le q-1),\\
 \Rcal(s,t)=\sum_{r=0}^q c_r\log\eta_r(s,t),\qquad
 \Scal(s,t)=\sum_{r=0}^q c_r\log\omega_{q+r}(s,t).
\end{gathered}
\end{equation}
Substitute (CPT29) at each of the four cutoffs.
The source prefix sums give coefficient \(-1\) at \(q,2q\)
and \(-2\) at \(q+1,\ldots,2q-1\).
The relation prefix sums give coefficient \(+1\) at \(0,q\)
and \(+2\) at \(1,\ldots,q-1\).  Thus
\begin{equation}\tag{CPT31}
 \boxed{\Bcal(s,t)=\Rcal(s,t)-\Scal(s,t),\qquad
 \sum_{r=0}^q c_r=2q.}
\end{equation}
This independently proves the signs and endpoint weights of
ACW20/SXR10 \cite{centralProvider,sourceRelation}.

If a source measure is multiplied by \(A>0\), all these minima
and every quotient metric are multiplied by \(A\).
The quotient determinant at each cutoff is multiplied by
\(A^q\), whose four signed logarithms cancel.
Equivalently each ratio \(\eta_r/\omega_{q+r}\) is unchanged.
Therefore the original arithmetic correction is exactly
\begin{equation}\tag{CPT32}
 \delta_k^\sigma
 :=\Bcal[m_k(y)dy]-\Bcal[\sigma(y)dy]
 =\Bcal(1,1)-\Bcal(0,0).
\end{equation}
This is the proved scalar cancellation of the original \(M^k\).

Use \(J=0,D=2q+1\), \(g=1\) in (CPT9), and abbreviate
\(\Lambda_0=\Lambda_{k,0,2q+1}\).
Every minimizing polynomial \(\chi u_r\), \(0\le r\le q\),
and each of its allowed coefficient variations belongs to
this complete relation space.
Differentiating \(\eta_r\) makes the coefficient-derivative
terms vanish by its defining orthogonality to
\(\chi\C[S]_{<r}\).  The remaining derivative is exactly
\[
 \partial_s\eta_r
 =\int_C H(y)|\chi(c+iy)u_{r;s,t}(c+iy)|^2
                                             d\nu_{s,t}(y).
\]
Equations (CPT11) and (CPT16) consequently prove
\begin{equation}\tag{CPT33}
 |\partial_s\log\eta_r(s,t)|\le\Lambda_0,\qquad
 |\partial_s\Rcal(s,t)|\le2q\Lambda_0.
\end{equation}
The common degree \(2q+1\) includes the output face needed when
one original linear factor multiplies a degree-\(2q\) relation.
For the monic relation band alone, \(D=2q\) also suffices and
gives its correspondingly sharper constant.

Splitting (CPT32) at the actual endpoint \((0,1)\), and using
(CPT31) on its central leg, gives the exact finite receiver
\begin{equation}\tag{CPT34}
\begin{split}
 \boxed{\delta_k^\sigma
 =\Bcal(0,1)-\Bcal(0,0)
 -[\Scal(1,1)-\Scal(0,1)]+E_{\rm rel}},\\
 E_{\rm rel}=\Rcal(1,1)-\Rcal(0,1),\qquad
 |E_{\rm rel}|\le2q\Lambda_0.
\end{split}
\end{equation}
For any function \(F\) on the square, define its oriented
rectangle difference
\(\Delta_\square F=F(1,1)-F(0,1)-F(1,0)+F(0,0)\).
Both central relation legs obey (CPT33), so
\begin{equation}\tag{CPT35}
 \boxed{\Delta_\square\Bcal+\Delta_\square\Scal
       =\Delta_\square\Rcal,\qquad
 |\Delta_\square\Rcal|\le4q\Lambda_0.}
\end{equation}
No claim of a small mixed derivative follows from this integrated
rectangle estimate; its derivation controls the two actual legs.

The sign estimate (CPT17) propagates through every monic
source minimum because
\(d\nu_{1,t}\le e^{K_h}d\nu_{0,t}\).
Thus
\(\Scal(1,t)-\Scal(0,t)\le2qK_h\).
Combining this with (CPT31),(CPT33) proves
\begin{equation}\tag{CPT36}
 \Bcal(1,t)-\Bcal(0,t)\ge-2qK_h-2q\Lambda_0
 \quad(0\le t\le1).
\end{equation}
At fixed packet (CPT13) implies \(q^A\Lambda_0\to0\) for every
fixed \(A>0\).  In particular the lower logarithmic-scale bound
\[
 \liminf_{k\to\infty}
 \frac{\Bcal(1,t)-\Bcal(0,t)}{q\log k}\ge0
\]
holds uniformly for \(t\in[0,1]\).
It leaves the source-band value and the exterior leg in
(CPT34) explicitly present.

\section{Exact mixed source-band kernel and complete cross terms}

Write \(\Delta_N(s,t)=\det H_N(s,t)\).
The monic source-product identity above gives
\begin{equation}\tag{CPT37}
 \Scal=
 \log\Delta_{2q-1}+\log\Delta_{2q}
          -\log\Delta_{q-1}-\log\Delta_q.
\end{equation}
Let \(f_N(y)=(1,c+iy,\ldots,(c+iy)^N)^T\).
With the conjugate-linear-first Gram convention,
\[
 H_N=\int\overline{f_N(y)}f_N(y)^T\,d\nu_{s,t}(y).
\]
Define its reproducing kernel by the original coefficient formula
\begin{equation}\tag{CPT38}
 K_N^{s,t}(y,z)
   =f_N(y)^T H_N(s,t)^{-1}\overline{f_N(z)}.
\end{equation}
If \(F(y)=f_N(y)^Ta\), matrix multiplication gives
\[
 \int K_N^{s,t}(y,z)F(z)d\nu_{s,t}(z)=F(y).
\]
Also \(K_N^{s,t}(z,y)=\overline{K_N^{s,t}(y,z)}\).
This verifies the kernel and its conjugations directly in the
original complex coefficient frame.

At each real point \(H_CH_O=0\).  Hence
\(\partial_s\partial_tH_N=0\), while
\[
 H_{N,C}=\int_C H(y)\overline{f_N(y)}f_N(y)^T\,d\nu_{s,t}(y),
 \quad
 H_{N,O}=\int_O H(z)\overline{f_N(z)}f_N(z)^T\,d\nu_{s,t}(z).
\]
Differentiating \(H_N^{-1}H_N=I\) gives
\(\partial_tH_N^{-1}=-H_N^{-1}H_{N,O}H_N^{-1}\).
The determinant derivative follows by differentiating the
finite determinant polynomial and using its adjugate:
\(\partial_s\log\det H_N=\Tr(H_N^{-1}H_{N,C})\).
Consequently
\begin{equation}\tag{CPT39}
\begin{split}
 \partial_s\partial_t\log\Delta_N
 &=-\Tr(H_N^{-1}H_{N,C}H_N^{-1}H_{N,O})\\
 &=-\int_C\int_O H(y)H(z)|K_N^{s,t}(y,z)|^2
                          d\nu_{s,t}(y)d\nu_{s,t}(z).
\end{split}
\end{equation}
For the second equality, expand the two integrals and use the
trace of the rank-one products.  The resulting scalar product
is \(K_N(y,z)K_N(z,y)=|K_N(y,z)|^2\).
Every integral is absolutely convergent by the moment bounds
proved after (CPT8); the finite inverse Gram introduces only
finite coefficients on the compact parameter square.

Combining all four terms of (CPT37) yields
\begin{equation}\tag{CPT40}
\begin{split}
 \partial_s\partial_t\Scal(s,t)
 =-\int_C\int_O H(y)H(z)\big[
 &|K_{2q-1}^{s,t}(y,z)|^2+|K_{2q}^{s,t}(y,z)|^2\\
 &-|K_{q-1}^{s,t}(y,z)|^2-|K_q^{s,t}(y,z)|^2
 \big]\,d\nu_{s,t}(y)d\nu_{s,t}(z).
\end{split}
\end{equation}
This is an exact expression for the source-band interaction,
with its original two consecutive windows.  Integrating over
the parameter square gives \(\Delta_\square\Scal\).
Equations (CPT35),(CPT40) therefore return the full signed
rectangle to this explicit mixed kernel integral with the
proved remainder \(4q\Lambda_0\).
Disjoint scalar supports made \(\partial_s\partial_tH_N=0\);
they do not make the two noncommuting compressed matrices
\(H_{N,C},H_{N,O}\) have zero product.

The relation determinant has the same exact formula.
If \(B_N\) is the full coefficient matrix of
\(\chi,\chi S,\ldots,\chi S^{N-q}\), its Gram and derivatives
are \(B_N^*H_NB_N,B_N^*H_{N,C}B_N,B_N^*H_{N,O}B_N\).
The relation version of (CPT39) therefore contains its complete
cross matrix.  Subtracting it from the source version in
(CPT29) gives the exact quotient mixed derivative.  At
\(N=q-1\) the relation matrix is empty, its determinant is one,
and its derivative is zero.  This proves the relation-to-source
and source-to-quotient maps without omitting the zero-cutoff
endpoint.

\section{Propagation into the earlier central statements}

Equations (CPT10)--(CPT11) reproduce the root-sensitive finite
annihilator localization C3--C7 from the supplied complete
proofs \cite{complete}, now uniformly on the entire square.
For the replacement \(\nu_{1,t}\to\nu_{0,t}\), the earlier
endpoint estimate multiplied that localization by
\(e^{L_k}-1\) and required a final smallness guard.
The proved derivative argument replaces its use by
\[
 \Lambda=L_k\varepsilon_{k,J,D},\quad
 e^{-\Lambda}G_{0,t}\preceq G_{1,t}\preceq e^\Lambda G_{0,t},
 \quad
 \norm{P_1-P_0}_{0,t}\le(e^\Lambda-1)\norm{P_0}_{0,t}.
\]
These are valid finite inequalities for every admitted
\(k\) satisfying the displayed density, degree and root guards.
The old endpoint inequalities remain preserved statements
on their own domains.

The previous central-leg relation radius in C17 is replaced,
for this exact interpolation, by \(2qL_k\varepsilon_{k,0,2q+1}\).
The uniform one-sided conclusion C21--C22 now has the radius
in (CPT36) for every exterior parameter.
The additional rectangle identity is (CPT35), with the
literal mixed kernel (CPT40).  These substitutions propagate
new constants and new path information; they do not evaluate
an uncomputed source-band integral.

For each localized original submodule and its actual fixed
subquotients, (CPT27) proves the quantitative correspondence.
For the complementary quotient, (CPT28) records its precise
Schur metric.  For the full signed determinant, (CPT34) and
(CPT40) retain the actual source-band and exterior quantities.
Thus no seven absolute allocation, proper-source
standard-window value, individual projected sign, or
Riemann-hypothesis conclusion is supplied by a substituted
rank or a two-class comparison in this note.

\appendix
\section{The retained local-density constants and proof}

This appendix states and verifies the particular central-density
provider used in (CPT14)--(CPT16), with the exact original constants
of ACW4--11 \cite{centralProvider}.  It uses the two original
source envelopes recorded below (CPT3), rather than any new
analytic hypothesis.
Put \(\eta=1/4\) and \(C_h^{\rm tilt}=C_{\rm tilt}\).
For \(|\theta|\le\alpha\), define the auxiliary density
\[
 p_\theta(y)=e^{\theta y}w_h(y)/M_h(\theta),\qquad
 \mu(\theta)=\ell_h'(\theta).
\]
Its unit mass is used only in this Fourier computation.  The
exact convolution identity is
\begin{equation}\tag{CPT41}
 m_k(x)=M_h(\theta)^k e^{-\theta x}p_\theta^{*k}(x).
\end{equation}
It follows by pulling the convolution back to its \(k\)
original variables, whose sum is \(x\) with Jacobian one.
Thus the original arithmetic mass is explicitly restored.

The envelope and \(M_h(\theta)\ge M_0\) give
\[
 V_*=\frac{2C_{\rm tilt}}{M_0}\operatorname B(3,p-3),
 \qquad
 H_*=
 2^{1+\eta}\left[
 \frac{2C_{\rm tilt}}{M_0}\operatorname B(3+\eta,p-3-\eta)
                      +\mu^{2+\eta}\right].
\]
Integrating \(y^2(1+|y|)^{-p}\), and using
\(|x-y|^{2+\eta}\le2^{1+\eta}(|x|^{2+\eta}+|y|^{2+\eta})\),
proves \(V(\theta)\le V_*\) and
\(\int|y-\mu(\theta)|^{2+\eta}p_\theta(y)dy\le H_*\).
The beta integrals converge, since \(p\ge7/2>3+\eta\).

On \([-1,1]\), (CPT41) at \(k=3\) and the original lower
convolution envelope give
\[
 p_\theta^{*3}(x)\ge
 c_h e^{-2\alpha}2^{-B_0}/M^3.
\]
Define
\[
 a_{\rm mix}=\min\{1/4,c_he^{-2\alpha}2^{-B_0}/M^3\},\quad
 \beta_{\rm mix}=2a_{\rm mix},\quad v_*=\beta_{\rm mix}/9.
\]
The triple convolution is the sum of \(\beta_{\rm mix}\)
times the uniform density on \([-1,1]\), and a nonnegative
residual of mass \(1-\beta_{\rm mix}\).
The variance of a mixture is at least the mixture-weighted
within-component variances, as follows by expanding each
centred square about the full mean.
It is therefore at least \(\beta_{\rm mix}/3\).
The variance of the three independent summands is
\(3V(\theta)\), proving \(V(\theta)\ge v_*>0\).
Continuity and differentiation of \(M_h\) through order two
follow from the same integrable envelope.
Consequently \(\mu(\theta)\) is strictly increasing and maps
\([-\alpha,\alpha]\) onto \([-\mu,\mu]\), establishing the
inverse used in (CPT14).

For the centred characteristic function
\(\phi_\theta(\zeta)=\int
 e^{i\zeta(y-\mu(\theta))}p_\theta(y)dy\),
Taylor's formula on \(|u|\le1\), and
\(2+|u|+u^2/2\le(7/2)|u|^{2+\eta}\) on \(|u|>1\), imply
\[
 |\phi_\theta(\zeta)-1+V(\theta)\zeta^2/2|
       \le4H_*|\zeta|^{2+\eta}.
\]
Put
\[
 \zeta_0=\min\{1,(2V_*)^{-1/2},
                  (v_*/(16H_*))^{1/\eta}\}.
\]
For \(|\zeta|\le\zeta_0\), the quadratic term is nonnegative
and the Taylor error is at most \(v_*\zeta^2/4\).  Hence
\(|\phi_\theta(\zeta)|\le e^{-v_*\zeta^2/4}\).
The triple-density mixture gives
\[
 |\phi_\theta(\zeta)|^3
 \le1-\beta_{\rm mix}
           +\beta_{\rm mix}|\sin\zeta/\zeta|.
\]
For \(|\zeta|\ge\zeta_0\) one has
\(|\sin\zeta/\zeta|\le1-\zeta_0^2/7\).
To check the constant, on \([\zeta_0,2]\) the positive
function \(\sin u/u\) decreases, and its Taylor upper bound
at \(0<\zeta_0\le1\) is
\(1-\zeta_0^2/6+\zeta_0^4/120
\le1-\zeta_0^2/7\).  Beyond \(2\), its absolute value is
at most \(1/2\), which is smaller.
Thus, with
\[
 r_*=(1-\beta_{\rm mix}\zeta_0^2/7)^{1/3}<1,\qquad
 C_2=\frac{4\pi(C_{\rm tilt}/M_0)^2}{2p-1},
\]
we have \(|\phi_\theta|\le r_*\) off the small interval.
The Fourier Plancherel identity \cite[(1.14.8)]{dlmfFourier}
in the convention
\(\widehat f(\zeta)=\int e^{i\zeta y}f(y)dy\) gives
\(\int|\phi_\theta|^2\le C_2\).
Relative to the DLMF convention
\(\mathcal F f=(2\pi)^{-1/2}\int e^{i\zeta y}f(y)dy\),
our transform is \(\sqrt{2\pi}\mathcal F f\).  Thus its
squared \(L^2\) norm is \(2\pi\int|f|^2\); integrating
the squared density envelope yields exactly \(C_2\).

Let
\[
 C_0=4H_*+V_*^2/8,\qquad a_k^{\rm loc}=v_*(k-1)/4,
\]
and define the finite error
\begin{equation}\tag{CPT42}
\begin{split}
 A_k^{\rm loc}=\frac1{2\pi}\bigg[
 &C_0k\,\Gamma((3+\eta)/2)
                    (a_k^{\rm loc})^{-(3+\eta)/2}
 +C_2r_*^{\,k-2}\\
 &+\frac{2e^{-kv_*\zeta_0^2/2}}{kv_*\zeta_0}
 \bigg],\qquad
 \epsilon_k^{\rm loc}=\sqrt{2\pi kV_*}\,A_k^{\rm loc}.
\end{split}
\end{equation}
On the small interval,
\[
 |\phi_\theta(\zeta)-e^{-V(\theta)\zeta^2/2}|
 \le C_0|\zeta|^{2+\eta}.
\]
This follows from the previous Taylor bound and the exponential
remainder \(V_*^2\zeta^4/8\), using
\(|\zeta|^4\le|\zeta|^{2+\eta}\) when \(|\zeta|\le1\).
Factor the difference of the \(k\)-th powers as a sum of \(k\)
terms.  Both factors have modulus at most
\(e^{-v_*\zeta^2/4}\), so its integral is bounded by
\[
 C_0k\int_\R|\zeta|^{2+\eta}
             e^{-a_k^{\rm loc}\zeta^2}\,d\zeta
 =C_0k\,\Gamma((3+\eta)/2)
                      (a_k^{\rm loc})^{-(3+\eta)/2}.
\]
The characteristic-function integral on the complementary
region is at most \(C_2r_*^{k-2}\), by its \(L^2\) bound.
The corresponding Gaussian tail is at most
\(2e^{-kv_*\zeta_0^2/2}/(kv_*\zeta_0)\), by replacing
the factor \(|\zeta|/\zeta_0\) with its lower bound one
before integrating the Gaussian derivative.
Fourier inversion therefore yields
\begin{equation}\tag{CPT43}
 \left|p_\theta^{*k}(k\mu(\theta))
           -(2\pi kV(\theta))^{-1/2}\right|
 \le A_k^{\rm loc}.
\end{equation}
All bounds are uniform for the closed interval
\(|\theta|\le\alpha\), including its two endpoints.
The use of inversion is justified by
\(\phi_\theta^k\in L^1\), already implied by the \(L^2\)
bound and \(|\phi_\theta|\le1\) for \(k\ge2\).
Multiplying (CPT43) by \(\sqrt{2\pi kV(\theta)}\) proves
the relative error in (CPT14), with
\(\epsilon_k^{\rm loc}=O_h(k^{-1/8})\).
In particular its stated finite guard is eventually satisfied.

Finally apply (CPT41) with \(\theta=\theta(u)\),
insert the relative-error expression, and divide by
\(M^k\sigma(ku)\).  Writing the Gamma density as
\[
 \sigma(ku)=c_{\sigma}^{\rm env}(ku)
                 e^{-\alpha k|u|}(1+k|u|)^{-1/2}
\]
gives (CPT14) by exact logarithms.
This completes the density-to-contrast calculation used by
the new path estimates.

\begin{thebibliography}{9}
\bibitem{complete}
Split-Zero mathematical continuation,
\emph{Complete additive proofs}, 18 September 2026,
source file \texttt{COMPLETE\_PROOFS.md},
central arithmetic density note, equations C1--C22.
The new path transport in this note extends its C3--C13,
C17 and C21--C22; all original source bytes are preserved.

\bibitem{providers}
Split-Zero mathematical continuation,
\emph{The full tilted mass in the original arithmetic tail;
A global tilted-mass transfer with the central source retained},
source file \texttt{RTM\_GTM\_original\_providers.tex},
RTM1--RTM3 and GTM1, GTM7.
This is the original-source provider carried inside the
18 September 2026 signed arithmetic value package.

\bibitem{centralProvider}
Split-Zero mathematical continuation,
\emph{Central source window},
source file \texttt{CENTRAL\_SOURCE\_WINDOW.md},
ACW1--ACW13 and ACW20--ACW26,
carried inside the 18 September 2026 signed arithmetic
value package.  This note reconstructs its needed Fourier
estimate in the appendix.

\bibitem{sourceRelation}
Split-Zero mathematical continuation,
\emph{The signed return through the actual adjacent relation and quotient},
source file \texttt{SXR\_original\_source\_relation.tex},
SXR1--SXR10, carried inside the same package.
The finite Schur and signed-band identities are proved
again here with the two-parameter measure.

\bibitem{dlmfLegendreNorm}
T.\ H.\ Koornwinder, R.\ Wong, R.\ Koekoek, R.\ F.\ Swarttouw
and W.\ P.\ Reinhardt,
NIST Digital Library of Mathematical Functions,
Chapter 18, Table 18.3.1, Legendre orthogonality and leading
coefficient conventions.
\url{https://dlmf.nist.gov/18.3#T1}.
Consulted 19 September 2026.

\bibitem{dlmfLegendreRec}
T.\ H.\ Koornwinder, R.\ Wong, R.\ Koekoek, R.\ F.\ Swarttouw
and W.\ P.\ Reinhardt,
NIST Digital Library of Mathematical Functions,
Chapter 18, Table 18.9.1, Legendre recurrence.
\url{https://dlmf.nist.gov/18.9#T1}.
Consulted 19 September 2026.

\bibitem{dlmfMP}
T.\ H.\ Koornwinder, R.\ Wong, R.\ Koekoek, R.\ F.\ Swarttouw
and W.\ P.\ Reinhardt,
NIST Digital Library of Mathematical Functions,
Chapter 18, equation 18.22.8, Meixner--Pollaczek recurrence.
\url{https://dlmf.nist.gov/18.22.E8}.
Consulted 19 September 2026.

\bibitem{dlmfFourier}
NIST Digital Library of Mathematical Functions,
Section 1.14(i), especially (1.14.1), (1.14.4) and (1.14.8),
Fourier inversion and Parseval identities.
Its source notes cite E.\ C.\ Titchmarsh,
\emph{Introduction to the Theory of Fourier Integrals},
second edition, Chelsea, 1986, pp.\ 3--15, 42, 50--59.
\url{https://dlmf.nist.gov/1.14#i}.
Consulted 19 September 2026.
\end{thebibliography}
\end{document}
